\aufzaehlungfuenf{Pemetaan
\maabbeledisp {\pi_{\text{vert} }} {TE= TX \times \R \times \R } { p^*E  = X \times \R \times \R
} {(P,u,v,w)} { (P,u,  \omega(P,v) + w)
} {}
bersifat linear dalam
\mathkor {} {v} {dan} {w} {,}
untuk $(P,u)$ yang tetap, sehingga merupakan suatu homomorfisme bundel vektor. Ketergantungannya pada basis $E$ terdiferensialkan secara kontinu karena $\omega$ terdiferensialkan secara kontinu. Pemetaan ini merupakan suatu pemisahan bagi penyertaan alami
\maabb {} {p^*E} { TE
} {}
karena
\mathl{(P,u,w)}{} dipetakan ke
\mathl{(P,u,0,w)}{} dan kemudian ke

\mavergleichskette
{\vergleichskette
{ (P,u, \omega(P,0) + w)
}
{ = }{ (P,u,  w)
}
{  }{
}
{    }{
}
{   }{
}
}
{}{}{.}
}{Untuk
\maabb {f} { U } {\R
} {}
kita perlu meninjau komposisi

\mathdisp {TX \stackrel{T(f)}{  \longrightarrow } T E = TX \times \R \times \R \stackrel{ \pi_{\text{vert} } }{ \longrightarrow } X \times \R \longrightarrow \R} {  }

yang menghasilkan

\mathdisp {(P,v) \longmapsto  (P,f(P), v, (df)_P(v) ) \longmapsto (P, f(P), \omega(P,v) + (df)_P(v) ) \longmapsto  \omega(P,v) + (df)_P(v)} { . }

}{Suatu penampang horizontal ada jika dan hanya jika turunan vertikalnya lenyap; dalam hal ini, syarat tersebut berarti

\mavergleichskettedisp
{\vergleichskette
{ \omega(P,v) + (df)_P(v)
}
{ =} { 0
}
{ } {
}
{ } {
}
{ } {
}
}
{}{}{}
untuk semua

\mavergleichskette
{\vergleichskette
{ (P,v)
}
{ \in }{ TX
}
{  }{
}
{    }{
}
{   }{
}
}
{}{}{.}
}{Hal ini mengikuti dari kenyataan bahwa suatu bentuk eksak jika dan hanya jika secara lokal bentuk tersebut mempunyai suatu bentuk primitif.
}{Hal ini mengikuti dari hasil-hasil yang telah dibuktikan.
}
